\begin{Example}[From state labels to a drift score]
\label{ex:drift-score}
Consider the resource perspective of the service-desk log in \Cref{ex:service-desk}. Suppose the fitted model assigns $W_{\mathrm a},W_{\mathrm b},W_{\mathrm c}$ to state $\mathsf R$, in which Frontline resolves both cases. It assigns $W_{\mathrm d},W_{\mathrm e},W_{\mathrm f}$ to state $\mathsf S$, in which Specialist resolves the focal case. Choose a signal width twice the feature width, with $b(W_{\mathrm a})$ on a signal-window boundary. The grouping rule in \Cref{def:state-distributions} yields
\[
  V_{\mathrm A}=W_{\mathrm a}\cup W_{\mathrm b},\qquad
  V_{\mathrm B}=W_{\mathrm c}\cup W_{\mathrm d},\qquad
  V_{\mathrm C}=W_{\mathrm e}\cup W_{\mathrm f},
\]
as depicted in \autoref{fig:approach-signal-windows}.

The state distributions $P_{\mathrm A}$, $P_{\mathrm B}$, and $P_{\mathrm C}$ count feature-window labels within the respective signal windows and divide by the number of labels. All labels in $V_{\mathrm A}$ are $\mathsf R$, the labels in $V_{\mathrm B}$ are evenly split, and all labels in $V_{\mathrm C}$ are $\mathsf S$ (\autoref{fig:approach-state-distributions}). The middle window mixes the earlier and later patterns because it spans the change.

To score $V_{\mathrm C}$, the previous-window reference is $Q_{\mathrm C}^{\mathrm{prev}}=P_{\mathrm B}$. The recent-window reference that uses both earlier signal windows is $Q_{\mathrm C}^{\mathrm{recent}}=\operatorname{mean}(P_{\mathrm A},P_{\mathrm B})$. Its share of $\mathsf S$ is smaller than in $P_{\mathrm B}$ because $V_{\mathrm A}$ contains only $\mathsf R$. The whole-log reference also includes $P_{\mathrm C}$. Its mean is $P_{\mathrm B}$ here because the exclusive $\mathsf R$ and $\mathsf S$ distributions balance each other (\autoref{fig:approach-reference-distributions}).

For these states, total variation reduces to the absolute difference in the share of $\mathsf S$. Write $\alpha=P_{\mathrm B}(\mathsf S)$ for the middle window's share of this state. This gives equal previous-window scores at $b(V_{\mathrm B})$ and $b(V_{\mathrm C})$.
\[
  D_{\mathrm{TV}}(P_{\mathrm B},P_{\mathrm A})
  =D_{\mathrm{TV}}(P_{\mathrm C},P_{\mathrm B})
  =\alpha.
\]
The recent-window reference gives $V_{\mathrm C}$ a higher score because its reference contains less $\mathsf S$. The whole-log mean gives $V_{\mathrm B}$ a zero score and the same positive score to both outer windows. \autoref{fig:approach-drift-scores} places each value at the start of the window it describes. Local scores are undefined until the required earlier signal windows are available.
\end{Example}

\begin{figure*}[!t]
  \centering
  \begin{tikzpicture}[
  kairo/base,
  x=0.93cm,
  row label/.style={anchor=east, align=right, inner sep=0pt, text=KairoMuted},
  window edge/.style={draw=KairoLine, line width=0.45pt},
  signal edge/.style={draw=KairoMuted, line width=0.55pt}
]
  \node[row label] at (-0.21,1.66) {Signal window};
  \node[row label] at (-0.21,1.04) {Feature window};
  \node[row label] at (-0.21,0.45) {State at start};

  \foreach \fromX/\toX/\windowName/\shade in
    {0/4/A/KairoPaper,4/8/B/KairoTealTint,8/12/C/KairoPaper} {
    \fill[fill=\shade] (\fromX,1.41) rectangle (\toX,1.91);
    \draw[signal edge] (\fromX,1.41) rectangle (\toX,1.91);
    \node at ({(\fromX+\toX)/2},1.66) {$V_{\mathrm{\windowName}}$};
  }

  \foreach \fromX/\windowName/\state/\shade/\ink in
    {0/a/R/KairoBlueTint/KairoBlue,2/b/R/KairoBlueTint/KairoBlue,
     4/c/R/KairoBlueTint/KairoBlue,6/d/S/KairoAmberTint/KairoAmber,
     8/e/S/KairoAmberTint/KairoAmber,10/f/S/KairoAmberTint/KairoAmber} {
    \fill[fill=\shade] (\fromX,0.80) rectangle ++(2,0.48);
    \draw[window edge] (\fromX,0.80) rectangle ++(2,0.48);
    \node at ({\fromX+1},1.04) {$W_{\mathrm{\windowName}}$};
    \draw[window edge] (\fromX,0.80) -- (\fromX,0.45);
    \fill[fill=\ink] (\fromX,0.45) circle[radius=1.8pt];
    \node[anchor=west, inner sep=2.5pt, text=\ink] at (\fromX,0.45)
      {$\mathsf{\state}$};
  }
  \foreach \boundary in {0,4,8,12} {
    \draw[signal edge] (\boundary,1.41) -- (\boundary,0.80);
  }
  \draw[KairoMuted, -{Latex[length=1.65mm]}, line width=0.45pt]
    (0,0.08) -- (12,0.08) node[above left, inner sep=1.5pt] {Time};
\end{tikzpicture}

  \caption{Signal windows in the service-desk example. In the middle signal window, Frontline resolves both cases in the earlier feature window and Specialist resolves the focal case in the later one.}
  \label{fig:approach-signal-windows}
\end{figure*}

\begin{figure*}[!t]
  \centering
  \begin{tikzpicture}[kairo/base, x=1cm, y=1cm]
  \foreach \fromX/\windowName/\barDescription in {
    0/A/{Frontline only},
    4.50/B/{Equal shares},
    9.00/C/{Specialist resolves the focal case}%
  } {
    \node[anchor=west, inner sep=0pt] at (\fromX,1.32) {$P_{\mathrm{\windowName}}$};
    \node[anchor=east, inner sep=0pt, text=KairoMuted] at ({\fromX+4.10},1.32)
      {Window $V_{\mathrm{\windowName}}$};
    \node[anchor=north, inner sep=0pt, text=KairoMuted] at ({\fromX+2.05},0.53)
      {\barDescription};
  }

  \foreach \fromX/\toX/\state/\shade/\ink in {
    0/4.10/R/KairoBlueTint/KairoBlue,
    4.50/6.55/R/KairoBlueTint/KairoBlue,
    6.55/8.60/S/KairoAmberTint/KairoAmber,
    9.00/13.10/S/KairoAmberTint/KairoAmber%
  } {
    \fill[fill=\shade] (\fromX,0.65) rectangle (\toX,1.10);
    \draw[draw=\ink, line width=0.5pt] (\fromX,0.65) rectangle (\toX,1.10);
    \node[text=\ink] at ({(\fromX+\toX)/2},0.875) {$\mathsf{\state}$};
  }
\end{tikzpicture}

  \caption{Observed state distributions in the signal windows. The first and last windows contain only $\mathsf R$ and $\mathsf S$, respectively. The middle window contains equal shares.}
  \label{fig:approach-state-distributions}
\end{figure*}

\begin{figure*}[!t]
  \centering
  \begin{tikzpicture}[kairo/base, x=1cm, y=1cm]
  \foreach \fromX/\barTitle/\barSource/\ink in {
    0/{Current window}/{$P_{\mathrm C}$}/KairoInk,
    3.52/{Previous window}/{$P_{\mathrm B}$}/KairoMuted,
    6.88/{Recent mean}/{$\operatorname{mean}(P_{\mathrm A},P_{\mathrm B})$}/KairoTeal,
    10.24/{Whole-log mean}/{Equals $P_{\mathrm B}$ here}/KairoMuted%
  } {
    \node[inner sep=0pt, text=\ink, font=\footnotesize\bfseries]
      at ({\fromX+1.435},1.32) {\barTitle};
    \node[anchor=north, inner sep=0pt, text=KairoMuted]
      at ({\fromX+1.435},0.53) {\barSource};
  }
  \draw[KairoLine, line width=0.5pt] (3.19,0.22) -- (3.19,1.45);

  \foreach \fromX/\toX/\state/\shade/\ink in {
    0/2.87/S/KairoAmberTint/KairoAmber,
    3.52/4.955/R/KairoBlueTint/KairoBlue,
    4.955/6.39/S/KairoAmberTint/KairoAmber,
    6.88/9.0325/R/KairoBlueTint/KairoBlue,
    9.0325/9.75/S/KairoAmberTint/KairoAmber,
    10.24/11.675/R/KairoBlueTint/KairoBlue,
    11.675/13.11/S/KairoAmberTint/KairoAmber%
  } {
    \fill[fill=\shade] (\fromX,0.65) rectangle (\toX,1.10);
    \draw[draw=\ink, line width=0.5pt] (\fromX,0.65) rectangle (\toX,1.10);
    \node[text=\ink] at ({(\fromX+\toX)/2},0.875) {$\mathsf{\state}$};
  }
\end{tikzpicture}

  \caption{Reference distributions for scoring $V_{\mathrm C}$. The recent-window reference has less $\mathsf S$ than the previous-window and whole-log references in this example.}
  \label{fig:approach-reference-distributions}
\end{figure*}

\begin{figure*}[!t]
  \centering
  \begin{tikzpicture}[
  kairo/base,
  x=1cm,
  y=1cm,
  grid line/.style={draw=KairoLine, line width=0.45pt},
  value line/.style={line width=0.91pt, line cap=round}
]
  \node[anchor=east, inner sep=0pt, text=KairoMuted] at (0.08,2.32)
    {Score};
  \foreach \level/\tickLabel in {0.38/$0$,1.10/$\alpha$,1.46/$\tfrac32\alpha$} {
    \node[anchor=east, inner sep=0pt, text=KairoMuted] at (0.08,\level) {\tickLabel};
  }

  \foreach \fromX/\panelTitle/\ink in {
    0.24/{Previous window}/KairoBlue,
    4.65/{Recent mean}/KairoTeal,
    9.06/{Whole-log mean}/KairoMuted%
  } {
    \node[font=\footnotesize\bfseries, inner sep=0pt, text=\ink]
      at ({\fromX+2.02},2.32) {\panelTitle};
    \fill[KairoPaper] (\fromX,0.38) rectangle ++(4.04,1.63);
    \foreach \level in {1.10,1.46}
      \draw[grid line] (\fromX,\level) -- ++(4.04,0);
    \draw[KairoMuted, line width=0.5pt] (\fromX,0.38) -- ++(4.04,0);
    \foreach \offset/\windowName in {0.60/A,2.02/B,3.44/C} {
      \draw[KairoMuted, line width=0.45pt]
        ({\fromX+\offset},0.38) -- ++(0,-0.06);
      \node[anchor=north, text=KairoMuted, inner sep=0pt]
        at ({\fromX+\offset},0.25) {$b(V_{\mathrm{\windowName}})$};
    }
  }

  \draw[value line, KairoBlue] (2.26,1.10) -- (3.68,1.10);
  \foreach \x in {2.26,3.68}
    \fill[KairoBlue] (\x,1.10) circle[radius=2.0pt];
  \fill[KairoTeal] (8.09,1.46) circle[radius=2.0pt];
  \draw[value line, KairoMuted] (9.66,1.10) -- (11.08,0.38) -- (12.50,1.10);
  \foreach \x/\y in {9.66/1.10,11.08/0.38,12.50/1.10}
    \fill[KairoMuted] (\x,\y) circle[radius=2.0pt];
\end{tikzpicture}

  \caption{Illustrative total-variation scores. The scale uses the middle window's share $\alpha$ of state $\mathsf S$. Missing points indicate that the reference needs more history. The whole-log reference gives the middle window a score of zero.}
  \label{fig:approach-drift-scores}
\end{figure*}
